% Generated draft: requires mathematical and source review.
\documentclass[11pt]{article}
\usepackage[margin=25mm]{geometry}
\usepackage{fontspec}
\setmainfont{Latin Modern Roman}
\usepackage{amsmath,amssymb,mathtools,unicode-math}
\setmathfont{Latin Modern Math}
\usepackage{enumitem,needspace,xurl,hyperref}
\setlength{\parindent}{0pt}
\setlength{\parskip}{5pt}
\setlength{\emergencystretch}{3em}
\newcommand{\N}{\mathbb{N}}
\newcommand{\Z}{\mathbb{Z}}
\newcommand{\Q}{\mathbb{Q}}
\newcommand{\R}{\mathbb{R}}
\newcommand{\C}{\mathbb{C}}
\newcommand{\F}{\mathbb{F}}
\newcommand{\A}{\mathbb{A}}
\newcommand{\PP}{\mathbb{P}}
\newcommand{\E}{\mathbb{E}}
\newcommand{\eps}{\varepsilon}
\newcommand{\dd}{\,\mathrm{d}}
\begin{document}
\section*{UnsolvedMath ID: 700724}
\textbf{Source problem code:} catalogue entry
\subsection{Definitions and mathematical statement}
A group \(H\) is \emph{polycyclic} if it has a finite subnormal series
\[
1=H_0\trianglelefteq H_1\trianglelefteq\cdots\trianglelefteq H_n=H
\]
whose factors \(H_{i+1}/H_i\) are cyclic; finite and infinite cyclic factors are allowed. For subgroups \(A,B\leq G\), write
\[
AB=\{ab:a\in A,\ b\in B\}.
\]
Thus \(G=AB\) requires every element of \(G\) to be such a product, and does not merely require \(A\) and \(B\) to generate \(G\).

The supplied record does not define the term \emph{radical group}. Let \(\mathcal R\) provisionally denote the class intended by that term in the source. The question of D. I. Zaitsev is whether
\[
\forall G\in\mathcal R\quad\forall A,B\leq G,\qquad
\bigl(G=AB\ \text{and}\ A,B\ \text{are polycyclic}\bigr)
\Longrightarrow G\ \text{is polycyclic}.
\]
The meaning of \(\mathcal R\) must be verified before this is a fully specified mathematical statement; no definition of radical is selected here.

The source additionally reports an affirmative answer for soluble groups and for hyperabelian groups, citing the two papers listed below. These are supplied background assertions, not independently verified results in this draft.
\subsection{Short English statement}
D. I. Zaitsev asks whether a radical group that is the setwise product of two polycyclic subgroups must itself be polycyclic. The supplied record reports affirmative answers for soluble and hyperabelian groups. Its meaning of ``radical'' remains unspecified and requires review.
\subsection{Sources}
\begin{itemize}
\item Local UnsolvedMath registry, supplied record with permanent ID \(700724\), exact source problem code \texttt{catalogue entry}, titled \emph{Kourovka Notebook Problem 11.31}. The registry labels the problem open; this label is supplied source data, not an independently verified literature assessment.
\item \emph{The Kourovka Notebook (Unsolved Problems in Group Theory)}, original problem 11.31, 11th issue (1990), proposed by D. I. Zaitsev. The supplied record identifies the 21st issue edition (2026), \texttt{arXiv:1401.0300v48}. Its collection description names E. I. Khukhro and V. D. Mazurov as editors.
\item D. I. Zaitsev, \emph{Math. Notes}, \textbf{29}, no. 4 (1981), 247--252. Supplied citation accompanying the reported affirmative cases; no article title or link is supplied.
\item J. C. Lennox and J. E. Roseblade, \emph{Math. Z.}, \textbf{170} (1980), 153--154. Supplied citation accompanying the reported affirmative cases; no article title or link is supplied.
\end{itemize}
\end{document}
